\begin{enumerate}
  \item
  \begin{enumerate}
    \item
    \begin{enumerate}[start=1, label=Caso \arabic*:, leftmargin=*]
      \item Las tres formas son: $1+1+1$, $1+2$ y $2+1$.
      \item Las dos formas son: $3+2$ y $2+3$.
      \item La única forma es dar un paso de $4$.
    \end{enumerate}

    \item
    \[
    \exists\, a, b \in \mathbb{N} \cup \{0\} : ai + bj = n
    \]

    \item
    \[
    f(n) =
    \begin{cases}
    1 & \text{si } n = 0 \\
    0 & \text{si } n < 0 \\
    f(n - i) + f(n - j) & \text{si } n > 0
    \end{cases}
    \]
  \end{enumerate}

  \item 
  \leavevmode
  \begin{pseudocode}
def formas(n, i, j):
    if n == 0: return 1
    if n < 0: return 0
    return formas(n - i, i, j) + formas(n - j, i, j)
  \end{pseudocode}
  
  {\small Se reproduce la parte conceptualmente significativa de la implementación; el archivo \texttt{solucion\_tarea\_1.py} contiene la implementación completa.}
\end{enumerate}
